\begin{table}[t!]
\caption{Characteristics of featured and non-featured playlists.\label{tab:promoplaylist}}
\smallskip{}
\raggedright{}
\resizebox{\textwidth}{!}{
\begin{tabular}{l S[table-align-text-post=false] S[table-align-text-post=false] }
\hline 
\multicolumn{1}{l}{Variable} & \multicolumn{1}{c}{Featured playlists} & \multicolumn{1}{c}{Non-featured playlists} \\ \hline 
\#\#\# & \#2\# & \#2\# \tabularnewline
\#\#\# & \#2\# & \#2\# \tabularnewline
\#\#\# & \#2\# & \#2\# \tabularnewline
\#\#\# & \#2\# & \#2\# \tabularnewline
\#\#\# & \#2\# & \#2\# \tabularnewline
\#\#\# & \#2\# & \#2\# \tabularnewline
\#\#\# & \#2\# & \#2\# \tabularnewline
\hline  

\end{tabular}
}
\begin{flushleft}
\vspace{-0.1cm}
\footnotesize{Notes: Reported across 34,483 playlists in our main sample, shown for playlists are featured at least once (vs. never featured). The table reports the number of playlists, the playlists' average followers, the share of playlists owned and curated by Spotify (in \%), the average share of major label content available on a playlist (in \%), the average number of tracks available on a list in a week, the likelihood that a playlist is updated in a given week, and the average track length (in minutes).}
\end{flushleft}
\end{table}